\chapter{Полнографовое вложение выровненного изотипического parent}
\label{ch:v8-baryon-c0-isotypic-channel-full-graph-aligned-parent}

\section{Коррекция carrier-счёта}

Предыдущий гейт правильно установил отсутствие полного edge-селектора, но
перенёс число нежелательных координат разреженной целевой опоры на число
полей вне $B\oplus M$. Эти числа различаются. Девять новых строгих рёбер
распадаются как
\begin{equation}
 E_{\rm new}=E_B^{\rm new}\sqcup E_M^{\rm new}\sqcup E_z,
 \qquad |E_{\rm new}|=9,quad (|E_B^{\rm new}|,|E_M^{\rm new}|,|E_z|)=(2,4,3).
 \label{eq:v8-full-graph-aligned-edge-decomposition}
\end{equation}
Здесь
\begin{align}
 E_B^{\rm new}&=\{L_LZ_R,Y_LZ_R\},\nonumber\\
 E_M^{\rm new}&=\{X_LZ_R,Z_Le_R,Z_LX_R,Z_LZ_R\},\nonumber\\
 E_z&=\{Z_LY_R,Z_Lu_R,Z_Ld_R\}.
 \label{eq:v8-full-graph-aligned-edge-sets}
\end{align}
Следовательно, матрицы $B$ и $M$ уже содержат шесть новых рёбер. Из шести
блоков, не входящих в разреженную желаемую опору, три являются компонентами
$M$, и только три действительно отсутствуют в subparent.

Корректное плоское продолжение предыдущего функционала поэтому содержит
\begin{equation}
 z=(z_Y,z_u,z_d)\in\mathbb C^3,
 \qquad
 N_{\rm omitted}^{\mathbb R}=6,
 \label{eq:v8-full-graph-aligned-correct-flat-count}
\end{equation}
а не шесть комплексных полей и двенадцать вещественных нулей.

\section{Условное положительное полнографовое завершение}

Минимальная квадратичная форма, включающая все девять новых рёбер, равна
\begin{equation}
 \mathcal S_{\rm full}
 =\mathcal S_{BM}
 +\mu_Y|z_Y|^2+\mu_u|z_u|^2+\mu_d|z_d|^2,
 \qquad
 \mu_Y,\mu_u,\mu_d>0.
 \label{eq:v8-full-graph-aligned-action}
\end{equation}
Каждое добавленное слагаемое калибровочно инвариантно и не меняет нулевое
множество выровненного сектора:
\begin{equation}
 z_Y=z_u=z_d=0,
 \qquad
 P_B+P_M=I_3,
 \qquad
 \ker M=\operatorname{im}P_B.
 \label{eq:v8-full-graph-aligned-zero-set}
\end{equation}
Квадратичный гессиан трёх новых комплексных полей имеет вид
\begin{equation}
 H_z=2\operatorname{diag}
 (\mu_Y,\mu_Y,\mu_u,\mu_u,\mu_d,\mu_d),
 \qquad
 \det H_z=64\mu_Y^2\mu_u^2\mu_d^2>0.
 \label{eq:v8-full-graph-aligned-extra-hessian}
\end{equation}
В нормировочной точке $\mu_Y=\mu_u=\mu_d=1$ полный восемнадцатимерный
вещественный срез имеет
\begin{equation}
 \Spec\Hess\mathcal S_{\rm full}
 =\{0\}^{\times4}\cup\{2\}^{\times6}
 \cup\{8\}^{\times5}\cup\{24\}\cup\{26\}^{\times2},
 \label{eq:v8-full-graph-aligned-spectrum}
\end{equation}
то есть
\begin{equation}
 (n_-,n_0,n_+)=(0,4,14).
 \label{eq:v8-full-graph-aligned-signature}
\end{equation}
Все неорбитальные плоские моды условно подняты.

\section{Три независимых массовых коэффициента}

Три дополнительных блока имеют различные endpoint-типы:
\begin{equation}
 \begin{aligned}
 Z_LY_R&:\text{ слабый дублет},\\
 Z_Lu_R&:\text{ верхний цветной канал},\\
 Z_Ld_R&:\text{ нижний цветной канал}.
 \end{aligned}
 \label{eq:v8-full-graph-aligned-extra-types}
\end{equation}
Поэтому калибровочная ковариантность не смешивает их и оставляет открытый
конус
\begin{equation}
 \mathcal C_\mu=(\mathbb R_{>0})^3,
 \qquad
 \dim\mathbb P\mathcal C_\mu=2.
 \label{eq:v8-full-graph-aligned-mass-cone}
\end{equation}
Равенство трёх масс является допустимой точкой, но не следствием симметрии.
Norm-square функционалы на пространствах представлений действительно дают
положительные страты и управляемые градиентные потоки
\cite{v8-full-graph-aligned-quiver}; спектральные расширения могут порождать
новые скалярные связи \cite{v8-full-graph-aligned-scalar}. Эти общие факты
не выбирают здесь конкретную точку $\mathcal C_\mu$.

\section{Что полное включение не выбирает}

Даже после добавления $z$ три нежелательных относительно разреженного меню
компоненты
\begin{equation}
 \{X_LZ_R,Z_Le_R,Z_LX_R\}\subset E_M^{\rm new}
 \label{eq:v8-full-graph-aligned-present-unwanted}
\end{equation}
остаются частью динамической матрицы $M$. Условие частичной изометрии
обычно делает её плотной на выбранном базисе, а общая правая орбита не
предпочитает координатную опору. Поэтому полнографовое включение снимает
плоские spectator-моды, но не восстанавливает старый разреженный рисунок.

Реестры принимают вид
\begin{equation}
 \begin{aligned}
 N_{\rm carrier\ embedding}&=7/7,
 &N_{\rm sparse\ selector}&=0/3,\\
 N_{\rm mass\ origin}&=0/4,
 &\mathcal M_{\rm new}&=\{\mu_Y,\mu_u,\mu_d,\text{ spectral origin}\}.
 \end{aligned}
 \label{eq:v8-full-graph-aligned-ledgers}
\end{equation}

\begin{theorem}[Корректное полнографовое завершение]
Матрицы $B$ и $M$ покрывают шесть из девяти новых строгих рёбер. Добавление
трёх оставшихся комплексных полей с положительными массами даёт полный
carrier и поднимает все неорбитальные плоские моды: на нормированном
вещественном срезе сигнатура равна $(0,4,14)$. Предыдущее число двенадцать
для дополнительных вещественных нулей заменяется точным числом шесть до
стабилизации.
\end{theorem}

\begin{proposition}[Свобода стабилизирующего конуса]
Калибровочные типы трёх добавленных полей неэквивалентны, поэтому симметрия
разрешает независимые положительные коэффициенты. Их общий масштаб и две
относительные координаты не выводятся из полнографового допуска. Кроме того,
полный carrier не выбирает разреженную coordinate-опору внутри $M$.
\end{proposition}

\begin{nogocriterion}[Запрет равных масс по удобству]
Нельзя объявлять $\mu_Y=\mu_u=\mu_d$ следствием общего следа до построения
единой суперсвязности и вычисления её кратностей. Нельзя также считать
плотную rank-two матрицу $M$ селектором трёх заранее названных рёбер.
\end{nogocriterion}

\begin{protocol}\begin{enumerate}
\item новых строгих рёбер: \textbf{девять};
\item новых рёбер внутри $B$: \textbf{два};
\item новых рёбер внутри $M$: \textbf{четыре};
\item действительно внешних полей: \textbf{три комплексных};
\item прежний счёт шести внешних полей: \textbf{скорректирован};
\item вещественных плоских spectator-мод до ремонта: \textbf{шесть};
\item положительное полнографовое завершение: \textbf{существует условно};
\item нормированная сигнатура: \textbf{$(0,4,14)$};
\item carrier-embedding ledger: \textbf{$7/7$};
\item разреженная coordinate-опора выбрана: \textbf{нет};
\item стабилизирующий конус: \textbf{$(\mathbb R_{>0})^3$};
\item относительных массовых свобод: \textbf{две};
\item mass-origin ledger: \textbf{$0/4$};
\item portal parent-origin: \textbf{не закрыт};
\item следующий гейт: \textbf{происхождение масс трёх дополнительных рёбер}.
\end{enumerate}\end{protocol}

Машинный сертификат сохранён в
\path{s2t_v8_baryon_c0_singlet_triplet_central_gap_isotypic_channel_full_graph_aligned_parent_embedding_gate_results.json}.

\begin{thebibliography}{9}
\bibitem{v8-full-graph-aligned-quiver}
M.~Harada, G.~Wilkin,
\emph{Morse Theory of the Moment Map for Representations of Quivers},
arXiv:0807.4734.
\bibitem{v8-full-graph-aligned-scalar}
C.~A.~Stephan,
\emph{New Scalar Fields in Noncommutative Geometry},
arXiv:0901.4676.
\end{thebibliography}